\pdfinfoomitdate=1
\pdftrailerid{}
\documentclass[10pt]{article}
\usepackage[a4paper,margin=0.70in]{geometry}
\usepackage[T1]{fontenc}
\usepackage{lmodern}
\usepackage{amsmath,amssymb,booktabs,microtype}
\title{Clean Fixed Circles Obstruct Isolated-Orbit\\
Determinants on a Heisenberg Nilmanifold}
\author{Route-A structural certificate C146}
\date{}
\begin{document}
\maketitle

\begin{abstract}
We freeze an explicit integer-lattice automorphism of the three-dimensional
Heisenberg group.  Its horizontal action is hyperbolic, while its center is
fixed pointwise.  Every iterate therefore owns a clean central fixed circle,
has singular ordinary isolated-orbit stability factor, and has zero Lefschetz
number.  The horizontal torus provides an exact isolated-point control through
iterate twenty.  These are source-side obstruction results, not a target or
arithmetic construction.
\end{abstract}

\section{The frozen lattice automorphism}
Write the real Heisenberg group as $H=\mathbb R^3$ with
\[
(x,y,z)(X,Y,Z)=(x+X,y+Y,z+Z+xY),
\qquad \Gamma=\mathbb Z^3.
\]
On the compact left quotient $N=\Gamma\backslash H$, freeze
\[
A=\begin{pmatrix}2&1\\1&1\end{pmatrix},\qquad
q(x,y)=x(x-1)+xy+\frac{y(y-1)}2,
\]
and
\begin{equation}
\Phi(x,y,z)=(2x+y,x+y,z+q(x,y)).                 \tag{1}
\end{equation}
The central correction is forced by the coordinate cocycle.  Direct expansion
gives
\[
q(v+w)-q(v)-q(w)=2xX+xY+Xy+yY,
\]
which equals $(2x+y)(X+Y)-xY$.  Thus (1) is a group homomorphism.
Moreover, $q(m,n)$ is integral for integer $m,n$, so the map preserves
$\Gamma$.  Since $A^{-1}$ is integral and the inverse central correction is
$-q(A^{-1}v)$, the inverse preserves $\Gamma$ as well.  Hence (1) is a lattice
automorphism and descends to $N$.

\section{Clean fixation and singular stability}
The embedded central circle
\[
C=\{[0,0,z]:z\in\mathbb R/\mathbb Z\}
\]
is fixed pointwise by $\Phi$, hence by every positive iterate.  The horizontal
eigenvalues are $(3\pm\sqrt5)/2$, neither of whose positive powers is one.
Consequently the horizontal fixed classes at each iterate are discrete, and
$C$ is a one-dimensional fixed component.  Along $C$, the derivative has
blocks $\left(\begin{smallmatrix}A^n&0\\ \ell_n&1\end{smallmatrix}\right)$.
The kernel equation first gives $(I-A^n)\delta v=0$, hence $\delta v=0$,
while $\delta z$ is free.  Therefore $\ker(I-D\Phi^n)=T C$, which is the
clean condition.

The derivative is block lower triangular.  For every $n\geq1$, its diagonal
blocks are $A^n$ and $1$, so
\begin{equation}
\det(I-D\Phi^n)=\det(I-A^n)(1-1)=0.              \tag{2}
\end{equation}
Thus the usual stability denominator for an isolated periodic orbit is
singular at every period.  We do not introduce a clean-family regularization.

Let $\alpha=dx$, $\beta=dy$, and $\gamma=dz-x\,dy$, with
$d\gamma=-\alpha\wedge\beta$.  Heisenberg cohomology has dimensions
$(1,2,2,1)$: by the Nomizu theorem the left-invariant complex computes the
cohomology of this compact nilmanifold, and direct calculation gives bases
$1$; $[\alpha],[\beta]$;
$[\alpha\wedge\gamma],[\beta\wedge\gamma]$; and
$[\alpha\wedge\beta\wedge\gamma]$.  Direct differentiation gives
\[
\Phi^*\alpha=2\alpha+\beta,\qquad
\Phi^*\beta=\alpha+\beta,\qquad
\Phi^*\gamma=\gamma-\alpha-\tfrac12\beta.
\]
The extra horizontal wedges are multiples of
$\alpha\wedge\beta=-d\gamma$, hence exact.  On the bases
$[\alpha],[\beta]$ and
$[\alpha\wedge\gamma],[\beta\wedge\gamma]$, pullback consequently has the
same trace $\operatorname{tr}(A^n)$; it acts as the identity in degrees zero
and three.  Therefore
\begin{equation}
L(\Phi^n)=1-\operatorname{tr}(A^n)+
\operatorname{tr}(A^n)-1=0.                    \tag{3}
\end{equation}

\section{Horizontal control and exact validation}
On the horizontal torus, fixed classes form
$\ker(A^n-I:\mathbb T^2\to\mathbb T^2)$, so their number is
\begin{equation}
|\det(A^n-I)|=\operatorname{tr}(A^n)-2=L_{2n}-2, \tag{4}
\end{equation}
where $L_k$ is the Lucas sequence.  Here $A=Q^2$ for
$Q=\left(\begin{smallmatrix}1&1\\1&0\end{smallmatrix}\right)$, proving the
last equality.  Exact evidence records (4), (2), and (3) through $n=20$.

The toral number in (4) is not a nilmanifold component count.  At $n=2$ the
horizontal class $v=(1/5,2/5)$ satisfies $(A^2-I)v=(2,1)$, but
$q(v)+q(Av)=0$.  If a point $(v,z)$ were fixed on the left quotient, some
$(2,1,k)\in\Gamma$ would obey $\Phi^2(v,z)=(2,1,k)(v,z)$.  Its central
coordinate would require $0=k+2(2/5)$, impossible for integral $k$.  Thus this
base class does not lift to a fixed circle, and we assert no full component
count for $\operatorname{Fix}(\Phi^n)$.

The independent checker passes 687 assertions, the SymPy reconstruction
passes 87 checks, byte replay passes, and all 31 hostile receipts (30 repaired-
hash semantic mutations and one stale hash) are rejected.

\section{Route-A boundary}
The strict tuple is $(\texttt{A1\_FAIL},\texttt{A2\_FAIL},
\texttt{A3\_FAIL},\texttt{A4\_FORMAL\_HINT})$, overall
\texttt{ROUTE\_A\_EXPLORATORY}.  The block determinant of $D\Phi$ is one, so
the descended map preserves Haar volume.  It gives a natural Koopman unitary
$U_\Phi f=f\circ\Phi$ on all of $L^2(N)$, with
$U_\Phi^n f=f\circ\Phi^n$ and hence the same iterate clock, but no
bridge to the singular clean-family weights; it is only a formal hint.  We
claim no isolated primitive-orbit ledger,
target divisor, target functional equation or counting law, arithmetic/local
factor, Euler factor, root number, automorphy, Hilbert--P\'olya construction,
or Route-B authorization.  The scope literal is
\texttt{NO\_BAD\_EULER\_OR\_ROOT\_NUMBER}.
\end{document}
